\section{Related Work}
\label{background}

AI-for-EDA (AI4EDA) uses learned models to support circuit analysis and optimization, including power, performance, and area estimation, timing prediction, and routability assessment. Predictors trained on circuit descriptions and tool-generated measurements can provide early feedback before expensive downstream analysis. The development of these predictors has progressively expanded from fitting models to engineered features, through learning circuit representations, to automating model design and experimentation. These approaches remain complementary; the progression highlights how the scope of automation has grown and which modeling decisions still require human expertise.

\subsection{Feature Engineering and Traditional Machine Learning}
\label{sec:feature_based_modeling}

Early learning-based circuit models typically started from a clearly specified prediction problem and a set of features selected by domain experts. Designers summarized relevant properties, such as switching activity, resource counts, or timing statistics, and fitted statistical or machine-learning models to the resulting feature vectors. Regression-based RTL power modeling by Bogliolo et al.~\cite{bogliolo2000regression}, for example, used linear regression and nonparametric extensions to relate power dissipation to input and output activity. For HLS, Dai et al.~\cite{dai2018fast} extracted features from synthesis reports and compared predictive models for estimating post-implementation quality of results.

These methods established the value of inexpensive learned approximations to EDA analyses. Their effectiveness nevertheless depended strongly on the information retained by the feature extractor. Once a circuit had been reduced to a fixed feature vector, the learner could exploit relationships among those features but could not recover omitted structural dependencies. Improving the predictor therefore often required an engineer to reconsider the feature set using knowledge of both the target EDA analysis and the learning method.

\subsection{Deep Learning for Circuit Prediction}
\label{sec:deep_circuit_modeling}

Deep learning shifted part of this effort from manually specifying feature interactions to learning them through neural architectures. Fully connected networks can learn nonlinear mappings from circuit descriptors, while convolutional neural networks (CNNs) can exploit spatial regularity in layout-derived inputs. RouteNet~\cite{xie2018routenet} illustrates the latter approach: it represents placement information as image-like inputs and applies CNNs to predict overall routability and design-rule-checking hotspots. This formulation connects the spatial structure of the prediction problem with an architecture suited to processing it.

Learned representations expanded the range of relationships that a predictor could capture, but model development still required substantial manual choices about input encoding, network architecture, training data, and optimization. A layout image is useful for spatial effects, whereas circuit connectivity need not follow proximity on that image. Likewise, a fully connected network operating on aggregate descriptors does not explicitly expose relationships between individual circuit elements. These limitations motivated models that process circuit structure more directly, alongside continued use of CNNs for spatial prediction tasks.

\subsection{Graph Neural Networks and Circuit-Specific Model Design}
\label{sec:semantic_codesign}

GNNs provide a natural representation for irregular circuit structures by treating operations, gates, pins, or other design objects as nodes and their dependencies as edges~\cite{lopera2021survey}. Compared with flattening a design into a fixed vector or rasterizing it onto a grid, graph-based modeling can explicitly retain connectivity and propagate information along task-relevant relations. Representative work spans HLS program graphs, netlist timing graphs, and logic representations, showing that the graph abstraction and learning procedure must be adapted to the property being predicted.

The existing literature illustrates three coupled dimensions of this adaptation. \emph{Representation (L1)} determines which circuit objects, relations, and features are visible to the model. GNN-DSE~\cite{sohrabizadeh2022automated} and Balor~\cite{murphy2024balor} develop graph-based HLS predictors, with Balor combining custom source-code graphs and hierarchical GNNs. \emph{Computation (L2)} determines how information is propagated and combined. PowerGear~\cite{lin2022powergear} uses a heterogeneous edge-centric GNN for power estimation, while the timing-engine-inspired model of Guo et al.~\cite{guo2022timing} structures prediction around timing propagation. \emph{Learning signal (L3)} determines the objectives and supervision used to train the representation. Guo et al. use local edge-delay prediction as auxiliary supervision, and DeepGate2~\cite{shi2023deepgate2} uses differences between sampled gates' truth tables to teach circuit functionality. These dimensions overlap: an effective improvement may require coordinated changes to the graph, the computation, and the training objective.

Thus, graph structure makes circuit information more accessible to a learner, but it also introduces additional modeling decisions. Engineers must construct appropriate graphs and datasets, choose node and edge attributes, design propagation and readout mechanisms, and select training objectives that reflect the target analysis. Data generation and labeling further depend on EDA tools and design constraints. Successful model development consequently draws on both machine-learning expertise and an understanding of circuits and EDA, motivating automation beyond parameter fitting alone.

\subsection{Neural Architecture Search and Automated Model Selection}
\label{sec:automated_gnn_design}

Neural architecture search (NAS) reduces the effort of manually selecting neural components by optimizing their combinations within a specified design space. Chang et al.~\cite{chang2020automatic} apply NAS to routability predictor development, providing an EDA-specific example of automated architecture construction. For graph learning, GraphNAS~\cite{gao2019graphnas} and Auto-GNN~\cite{zhou2022auto} use reinforcement-learning-based search to select GNN architectures, while the systematic design-space study of You et al.~\cite{you2020design} examines architectural and training choices across graph-learning tasks. These efforts make it easier to obtain competitive configurations without manually testing each candidate.

The remaining boundary is the design space itself. Available operators, connectivity patterns, and tunable settings must be represented in the search procedure, which evaluates candidates against a prescribed performance objective. The amount of manual architecture exploration is reduced, but experts still define the circuit representation, permissible transformations, and often the learning objective. If validation errors reveal a missing circuit relation or an unsuitable propagation rule, a fixed-space search can address that limitation only when the relevant change is already expressible. Automating this further step requires interpreting the model's behavior and creating new modeling strategies, rather than only selecting among encoded alternatives.

\subsection{LLM Agents and Autonomous Circuit Modeling}
\label{sec:llm_agents}

LLM-based research systems extend automation to proposing changes, editing code, running experiments, and interpreting results. The AI Scientist~\cite{lu2024aiscientist} connects idea generation with implementation, experimentation, and scientific reporting. AutoResearch~\cite{karpathy2026autoresearch} focuses on iterative modification and evaluation of language-model training code. Arbor~\cite{jin2026arbor} organizes autonomous research around a persistent tree of proposals, artifacts, evidence, and accumulated insights. These systems demonstrate mechanisms for carrying experimental findings into subsequent attempts and reducing repeated human intervention in the development loop.

Circuit modeling benefits from these capabilities, but useful feedback must also reflect the EDA problem being learned. An aggregate validation score establishes whether a predictor improves under the selected metric; additional circuit-specific analysis can help explain where it fails and what to change. For example, timing errors can be examined by path depth or cell type to investigate deficiencies in information propagation, while resource-prediction errors can be related to operation types or synthesis transformations. Such diagnostics turn circuit knowledge into evidence that can guide changes to representations, architectures, or training objectives. They complement general observations about convergence, capacity, and optimization with information about the physical or logical relationships the model is intended to capture.

CircuitModelPilot builds on the emerging autonomous-research workflow with this circuit-modeling focus. It uses task descriptions, circuit knowledge, and experimental feedback to analyze the current model, propose a revised modeling strategy, generate its implementation, and test the result. A model-design search tree retains alternative candidates and their outcomes so that later iterations can extend useful directions and reconsider unsuccessful ones. Relative to fixed-space NAS, the framework can introduce implementations that were not enumerated beforehand; relative to general research or training agents, its emphasis is on using EDA-specific evidence to make modeling revisions more targeted. The goal is to reduce the combined AI and EDA engineering effort needed to develop task-specific predictors and advance toward autonomous circuit modeling.
